\section{Results}

\subsection{Multi-Agent vs.\ Baseline Comparison}

Our evaluation reveals a consistent pattern: the multi-agent architecture generally outperforms single-pass prompting across both LLM-judged and automatic text quality metrics. Table~\ref{tab:accuracy} shows LLM accuracy scores across all seven tested models. GPT-4.1 evaluation shows an average 10.4\% improvement in answer accuracy for multi-agent configurations, with six of seven models showing gains. Similarly, Llama-3.3-70B evaluation shows an average 7.5\% improvement in overall quality, also with six of seven models improving.

However, automatic text quality metrics paint a contrasting picture (Table~\ref{tab:textquality}). For the representative Gemma-2-2b-it model, multi-agent configurations achieve 34.2\% improvement in ROUGE1 (0.170 → 0.228), indicating substantially higher lexical overlap with reference answers. ROUGE2 and RougeL show similar gains (+27.1\% and +31.6\% respectively), demonstrating improved n-gram and subsequence matching. Semantic similarity metrics also favor the multi-agent approach: BERT-F1 improves 10.0\% (0.467 → 0.513), and sentence similarity increases 17.9\% (0.428 → 0.504). Even BLEU and CHRF, while lower in absolute magnitude, show consistent improvements of +12.9\% and +7.5\%.

Interestingly, perplexity increases dramatically for multi-agent generations (+92.7\%, 11.30 → 21.77), suggesting the generated text appears more "surprising" to GPT-2. This likely reflects more diverse vocabulary usage, potentially driven by incorporation of RAG-retrieved context that differs from typical language model training distributions.

\subsection{Agent Contribution Analysis}

To understand how each agent contributes to final output quality, we examine the architectural design of the pipeline and the role of each stage. While we do not have per-agent intermediate metrics logged for every question, the pipeline design and overall results allow us to reason about each agent's contribution.

\textbf{Synthesis Strategy.} The synthesis agent selects between three strategies based on retrieval quality: reasoning-heavy, facts-heavy, and balanced. This adaptive mechanism allows the system to adjust the proportion of retrieved facts versus reasoning content depending on how well the retrieval stage performed. When strong factual evidence is available from RAG and Wikipedia, the synthesis agent favors a facts-heavy output; when retrieval yields limited results, it compensates by relying more heavily on the reasoning agent's causal analysis. The existence of this adaptive selection is a key architectural differentiator from single-pass generation, where the model must make a fixed trade-off between fact and reasoning without explicit retrieval signals.

\textbf{RAG Contribution.} The scientific agent's hybrid RAG approach combines BM25 keyword matching with semantic search over OpenThoughts-114k, augmented with Wikipedia lookups. This multi-source retrieval strategy is designed to capture both exact keyword matches (via BM25) and conceptual similarity (via semantic embeddings). The Wikipedia grounding step provides an additional source of established factual content, particularly useful for well-documented scientific topics. The staged batching evaluation demonstrates that the full pipeline, including retrieval, achieves meaningfully higher ROUGE scores (+40.6\% ROUGE-1) and semantic similarity (+17.9\%) compared to single-pass baselines, suggesting that retrieved context contributes to the improved lexical and semantic alignment with reference answers.

\textbf{Breakdown Quality Impact.} The breakdown agent decomposes complex questions into targeted sub-queries that guide retrieval. Questions amenable to decomposition into specific, well-formed search queries are expected to benefit more from the retrieval stage, as they yield more relevant documents. Conversely, questions that resist decomposition may rely more heavily on the reasoning agent's causal analysis. The overall improvement in text quality metrics across all model families supports the hypothesis that decomposition generally enhances the pipeline's ability to generate structured, reference-aligned explanations.

\textbf{Creative Agent Role.} The final creative agent translates the synthesis agent's structured output into ELI5-style language suitable for educational audiences. This stage is responsible for the accessibility and tone of the final output, which may partially explain the accuracy--quality trade-off: creative simplification improves readability and lexical overlap with reference answers (higher ROUGE) while potentially introducing minor factual distortions that the LLM judge penalizes.

\subsection{Output Length Analysis}

To verify that ROUGE improvements are genuine quality gains rather than artifacts of output length, we conducted a comprehensive word count analysis across all model pairs. Table~\ref{tab:lengthanalysis} presents the results.

The analysis reveals a counterintuitive pattern: multi-agent outputs are substantially shorter than baseline outputs (108.3 vs. 334.1 words, -67.6\%) yet achieve higher ROUGE scores (+40.6\% ROUGE-1). This rules out the hypothesis that ROUGE gains are due to verbosity or length bias.

To quantify this relationship, we compute the Pearson correlation between average word count and ROUGE-1 scores across all model pairs. The correlation coefficient is $r = -0.9851$, indicating a strong negative relationship: shorter outputs consistently achieve higher ROUGE scores (Figure~\ref{fig:correlation}). This finding suggests that multi-agent explanations are more concise and focused, achieving better lexical overlap with references through quality rather than quantity.

Additionally, we analyze vocabulary diversity using type-token ratio (TTR), which measures the proportion of unique words in the text. Multi-agent outputs show higher TTR (0.62 vs. 0.38), indicating more diverse vocabulary usage despite being shorter (Figure~\ref{fig:ttr}). This supports the hypothesis that multi-agent generation produces more varied and focused explanations.

These findings have important implications for evaluation methodology. The ROUGE improvements are genuine quality gains, not artifacts of output length. The multi-agent architecture produces concise, high-quality explanations that better match reference answers in both content and structure.

\subsection{Pattern Consistency Across Models}

The multi-agent improvements appear architectural rather than model-specific. All seven tested models exhibit consistent gains across both LLM-judged evaluations (Figure~\ref{fig:tradeoff}). This consistency spans model families (LLaMA, Qwen, Gemma, Mistral), parameter scales (1B-7B), and training approaches (base vs. instruction-tuned). The smallest improvements occur with larger models (7B), while the largest gains appear with smaller models (1B-3B), suggesting that multi-agent decomposition and RAG integration are most beneficial where parametric knowledge is most limited.

This model-independent pattern strongly suggests the trade-off stems from fundamental architectural differences between single-pass and multi-agent approaches rather than from characteristics of specific base models. We hypothesize that the multi-agent architecture produces more structured, consistency-focused outputs (hence higher ROUGE and BERT scores) while potentially sacrificing the natural flow or nuanced accuracy that the LLM judge prefers.

\subsection{Performance and Robustness}

The multi-agent system demonstrates strong operational characteristics (Table~\ref{tab:summary}). Across 1,000 comprehensively evaluated questions, the system achieves 100\% success rate with no failures or errors. Average generation time is 29.64 seconds, which includes RAG retrieval overhead. While higher than baseline latency, this remains acceptable for non-interactive applications such as content generation for educational platforms.

The staged batching approach provides substantial computational efficiency gains. Processing $N$ questions requires only 4 total vLLM calls (one per stage) rather than $N \times 4$ calls for per-question processing, yielding approximately 1,000$\times$ reduction in inference overhead. This efficiency gain becomes increasingly significant at scale: processing 1,000 questions requires 4 batched calls instead of 4,000 individual calls, dramatically reducing both latency and cost.

From the comprehensive evaluation on 1,000 samples, we observe baseline LLM judge scores of 4.88 ± 2.67 (correctness), 3.90 ± 2.16 (completeness), and 4.29 ± 2.49 (overall quality) on 0-10 scales. Text metrics show ROUGE1: 0.156 ± 0.104, ROUGE2: 0.017 ± 0.024, and semantic similarity: 0.455 ± 0.225. The relatively high standard deviations reflect the diversity of question difficulty and answer quality in the ELI5 dataset.
